\documentclass[12pt]{report}
\usepackage[spanish]{babel}
\usepackage[utf8]{inputenc}
\usepackage{amsmath,amsthm,amsfonts,amssymb}
\usepackage{Sweave}
\usepackage{graphicx}
\usepackage{multirow}
\usepackage{hyperref}
\usepackage{anysize}
\marginsize{1.78cm}{1.65cm}{1.78cm}{1.78cm} 

\title{\Huge Universidad Nacional de Loja \\ 
Área de la Energía las Industrias y los Recursos Naturales no Renovables \\
Ingeniería en Sistemas \\}


\author{\includegraphics[width=9cm, height=7cm]{logo-unl.jpg} \\\\ \href{http://www.iralis.org/?q=node\%2F10\&paso=10\&letra=\&id=7221}{ECINF7231} \\\\ \texttt{Martha Cristina Suntaxi Sarango \\ mcsuntaxis@unl.edu.ec \\\\ }}

\begin{document}
\Sconcordance{concordance:Analisis.tex:Analisis.Rnw:%
1 65 1 1 2 41 0 2 2 8 0 1 2 3 1}

\maketitle{\textbf {PARTE B}
\\[19]
\textbf{\large Análisis de Datos}
\\[8]
\textbf{¿Es aplicable la ingeniería de software cuando se elaboran webapps? Si es así, ¿cómo puede modificarse para que asimile las características únicas de éstas?}
\\[5]
Si se aplica, pero para esto el enfoque de ingenieria de software debe centrarse en ciertos requisitos o atributos propios de las webapps como la concurrencia, el uso intensivo de redes, disponibilidad, rendimiento y otros que son atributos de webapps.
\\[10]
\textbf{Dataset}
\\
Titanic
\\\\
\textbf{Descripción}
\\
Este conjunto de datos proporciona información sobre el destino de los pasajeros en el primer viaje fatal del trasatlántico Titanic, que se resumen de acuerdo a la situación económica (clase), el sexo, la edad y la supervivencia.
\\\\
\textbf{Formato}
\\
Una matriz de 4-dimensional resultante de cruzar los datos tabulados de 2201 observaciones sobre 4 variables. Las variables y sus niveles son los siguientes:
\\
\begin{table}[htb]
\begin{center}
\begin{tabular}{|l|l|l|}
\hline
No & Nombre & Niveles \\
\hline \hline
1	& Clase	& Primero, segundo, tercero, Tripulación \\ \hline
2	& Sexo	& Hombre, Mujer \\ \hline
3	& Años	& Niño, Adulto \\ \hline
4	& Sobrevivieron	& No, si \\ \hline
\end{tabular}
\end{center}
\end{table}
\\\\
\textbf{Detalles}
\\
Estos datos fueron recogidos originalmente por la Junta Británica de Comercio en su investigación del hundimiento.
\\
El hundimiento del Titanic es un evento famoso, y nuevos libros siguen siendo publicado sobre el tema.
\\\\
\textbf{Fuente}
\\
Dawson, Robert J. MACG. (1995), El 'Episodio inusual' Datos Revisited. Diario de Estadísticas de Educación, 3. \url{Https://www.amstat.org/publications/jse/v3n3/datasets.dawson.html}
\\\\
\textbf{Mostrar el dataset}

\begin{Schunk}
\begin{Sinput}
> Titanic
\end{Sinput}
\begin{Soutput}
, , Age = Child, Survived = No

      Sex
Class  Male Female
  1st     0      0
  2nd     0      0
  3rd    35     17
  Crew    0      0

, , Age = Adult, Survived = No

      Sex
Class  Male Female
  1st   118      4
  2nd   154     13
  3rd   387     89
  Crew  670      3

, , Age = Child, Survived = Yes

      Sex
Class  Male Female
  1st     5      1
  2nd    11     13
  3rd    13     14
  Crew    0      0

, , Age = Adult, Survived = Yes

      Sex
Class  Male Female
  1st    57    140
  2nd    14     80
  3rd    75     76
  Crew  192     20
\end{Soutput}
\end{Schunk}
\textbf{¿Cuál es el número total de casos en el dataset?}
\begin{Schunk}
\begin{Soutput}
Number of cases in table: 2201 
Number of factors: 4 
Test for independence of all factors:
	Chisq = 1637.4, df = 25, p-value = 0
	Chi-squared approximation may be incorrect
\end{Soutput}
\end{Schunk}
\begin{center}
\href{http://creativecommons.org/licenses/by-nc/4.0/}{\includegraphics[width=4cm, height=1.5cm]{cc.png}}
\end{center}
\end{document}
